% R32 section 4: results, gate H, obc0 (input by R32_floquet_spectra.tex)
\section{Results: the sixteen rules with the Hadamard gate}
\label{sec:results}

\subsection{The organising principle: how many neighbour patterns fire the gate}

Count, for each rule, the number $k$ of the four neighbour patterns $(x_{i-1},x_{i+1})\in\{00,01,10,11\}$ on which the
Hadamard acts. For $k=2$, also record whether the firing set depends on only one neighbour (``one-sided'': $60,102,153,195$)
or on both (``XOR'': $150$ fires on $01,10$, and $105$ on $00,11$). This single combinatorial number sorts the 16 rules into
classes with distinct spectral signatures (Table~\ref{tab:summary}, Fig.~\ref{fig:overview}):

\begin{center}
\small
\begin{tabular}{llp{0.5\textwidth}}
\toprule
class & rules & spectral signature (obc0) \\
\midrule
$k=0$ & $204$ & $U=I$ (identity) \\
$k=1$ & $108,156,198,201$ & exponential fragmentation; Poisson inside the pooled blocks \\
$k=2$ one-sided & $60,102,153,195$ & a tower of gate-algebra blocks; $\langle r\rangle$ near Poisson, drifting up with $N$.
  \emph{On a ring: no tower, chaotic} (Sec.~\ref{sec:pbc}) \\
$k=2$ XOR & $105,150$ & arithmetic spectrum, $3^m$ distinct quasienergies (W150) \\
$k=3$ & $54,57,99,147$ & \textbf{chaotic (GOE)}, plus a fixed number of exact gate-singlet eigenstates.
  On a ring: chaotic, CUE or COE by symmetry \\
$k=4$ & $51$ & $U=H^{\otimes N}$, non-interacting \\
\bottomrule
\end{tabular}
\end{center}

\begin{table}[t]
\centering
\scriptsize
\IfFileExists{tab_r32_summary_H.tex}{\resizebox{\textwidth}{!}{\begin{tabular}{llrlrrrrrrrr}
\toprule
$k$ & $W$ & $N$ & sym. & sectors & blocks & $\langle r\rangle_\mathrm{none}$ & $\langle r\rangle_\mathrm{full}$ & \#$r$ & distinct $\theta/2^N$ & $\bar\eta_\mathrm{big}$ & singlets \\
\midrule
0 & 204 & 16 & R & 65536 & 65536 & -- & -- & 0 & 1.53e-05 & -- & 0 \\
1 & 108 & 19 & R & 55405 & 124345 & 0.386 & 0.392 & 71048 & 0.0303 & 0.55 & 52688 \\
1 & 156 & 19 & none & 10946 & 121370 & 0.412 & 0.421 & 5486 & 0.00408 & 0.36 & 37812 \\
1 & 198 & 19 & none & 10946 & 121370 & 0.412 & 0.421 & 5486 & 0.00408 & 0.35 & 37812 \\
1 & 201 & 19 & R & 31572 & 90053 & 0.383 & 0.389 & 97266 & 0.068 & 0.54 & 39872 \\
2 one-sided & 60 & 14 & none & 15 & 105 & 0.384 & 0.395 & 7732 & 0.492 & 0.68 & 24 \\
2 one-sided & 102 & 14 & none & 15 & 105 & 0.384 & 0.395 & 7732 & 0.492 & 0.68 & 24 \\
2 one-sided & 153 & 14 & none & 1 & 15 & 0.398 & 0.421 & 8092 & 0.992 & 0.81 & 2 \\
2 one-sided & 195 & 14 & none & 1 & 15 & 0.398 & 0.421 & 8092 & 0.992 & 0.81 & 2 \\
2 XOR & 105 & 15 & R & 9 & 17 & 0.491 & 0.498 & 1512 & 0.0186 & 0.81 & 7 \\
2 XOR & 150 & 15 & R & 9 & 23 & 0.404 & 0.389 & 5987 & 0.1 & 0.66 & 7 \\
3 & 54 & 14 & Omega & 2 & 3 & 0.428 & 0.523 & 8127 & 0.992 & 0.84 & 1 \\
3 & 57 & 14 & none & 1 & 3 & 0.531 & 0.531 & 8190 & 1 & 0.96 & 2 \\
3 & 99 & 14 & none & 1 & 3 & 0.531 & 0.531 & 8190 & 1 & 0.96 & 2 \\
3 & 147 & 14 & Omega & 1 & 5 & 0.422 & 0.527 & 8117 & 0.991 & 0.89 & 4 \\
4 & 51 & 14 & R & 1 & 16384 & -- & -- & 0 & 0.000122 & -- & 16384 \\
\bottomrule
\end{tabular}
}}{\fbox{table pending}}
\caption{Per-rule summary at each rule's largest computed $N$ (obc0, gate $H$), generated by
\fn{spectral.analysis.write\_table}. sym.: spatial symmetry resolved ($R$ for odd $N$; the half-step $\Omega$ for even $N$).
sectors/blocks: Krylov sectors, and fully resolved gate-algebra blocks. $\langle r\rangle_\mathrm{none}$: whole spectrum as one
list. $\langle r\rangle_\mathrm{full}$: pooled over fully resolved blocks with $\ge20$ distinct half-levels, from \#$r$ ratios.
distinct $\theta/2^N$: number of distinct quasienergies over the whole spectrum. $\bar\eta_\mathrm{big}$: mean
$S_2/S_2^\mathrm{rand}(\mathrm{block})$ in the largest block. singlets: one-dimensional gate-algebra blocks inside Krylov
sectors of dimension $\ge3$, i.e.\ exact eigenstates of every individual gate.}
\label{tab:summary}
\end{table}

\begin{figure}[t]
\centering
\ifgraphic{../../figures/r32/fig_r32_overview_H.pdf}{\textwidth}
\caption{(a) Gap ratio of every rule at its largest $N$: unresolved (whole spectrum, open circles) and fully resolved
(Sec.~\ref{sec:symres}, filled). Lines: Poisson, GOE ($\beta=1$, the class of $U$), GUE. (b) Fraction of distinct
quasienergies in the whole spectrum (log scale). Rules grouped by the number $k$ of firing neighbour patterns.}
\label{fig:overview}
\end{figure}

\subsection{$k=3$: chaos, with exact scars that do not grow}
With open boundaries, the four rules on which the gate fires on all but one pattern are the only chaotic ones. (On a
ring, the one-sided $k=2$ rules join them; Sec.~\ref{sec:pbc}.)
\begin{itemize}
\item \textbf{57 and 99} (mirror images; computed independently, identical to all digits): $\langle r\rangle=0.5311\pm0.0028$
  at $N=14$ (8190 ratios), on the GOE value $0.5307$ from $N=10$ on: $0.525, 0.523, 0.527, 0.535, 0.531$ for
  $N=10\ldots14$. $P(r)$ follows the $\beta=1$
  surmise over the whole range (Fig.~\ref{fig:rule57}). The eigenstates in the big block are thermal \emph{within that
  block}: $\bar\eta=0.95$.
\item \textbf{147}: $\langle r\rangle_\mathrm{full}=0.526$ ($N=13$) and $0.5275\pm0.0028$ ($N=14$), but only after resolving
  the reflection (odd $N$) or the half step $\Omega$ (even $N$). Unresolved, the same spectra give $0.424$ and $0.422$. This is the textbook symmetry-mixing effect,
  and here it is the difference between ``integrable'' and ``chaotic''.
\item \textbf{54}: $\langle r\rangle_\mathrm{full}$ rises steadily, $0.441, 0.456, 0.502, 0.510, 0.523$ for
  $N=10\ldots14$, towards GOE. Its frozen vacuum $|0\ldots0\rangle$ is a separate Krylov sector.
\end{itemize}

\paragraph{Gate singlets.} Inside the chaotic block, each of these rules carries one-dimensional blocks of the gate
algebra: states that are eigenvectors of \emph{every individual gate} $g_i$, hence of $U$ for any order and any number of
gate applications. Their number does not depend on $N$ ($N=6\ldots13$):
\begin{itemize}
\item two for 57 and 99, and four for 147, all \emph{product} states ($S=0$) at $\theta\in\{0,\pi\}$;
\item exactly one for 54, at $\theta=0$, which is \emph{entangled}. Its entanglement \emph{decreases} with $N$:
  $S_\mathrm{vN}=0.470, 0.430, 0.405, 0.371, 0.345, 0.315, 0.291, 0.265, 0.243$ for $N=6\ldots14$ ($S_2$ from 0.35 to
  0.13). It is an area-law state that approaches a product state as the chain grows. (An earlier draft of this report said
  ``0.29 at every $N$''; that generalised from the single value at $N=12$ and is withdrawn.)
\end{itemize}
These are exact many-body scars in the sense of Refs.~\cite{shiraishi2017,moudgalya2022review}: non-thermal eigenstates
embedded in a GOE spectrum. The embedding is exact here, by the gate algebra, not approximate. Because they are
eigenstates of every gate, they are also stationary under any reordering of the circuit, which a Hamiltonian scar cannot
offer.

\paragraph{N-stable outliers in 57/99.} Besides the singlets, 57 has a handful of anomalously weakly entangled eigenstates
in the big block ($\eta\approx0.2$): 4 at $N=11$ and $13$, and 10 at $N=12$. Their quasienergies are stable with $N$
($\theta/\pi=\pm0.6877\to\pm0.6879$ from $N=11$ to $13$, and $\pm0.3122$ at $N=12$ and $13$). They do not form an equally
spaced tower, and their number does not grow. This points to a localised structure (a single defect in a frozen
background) rather than a PXP-type scar band. In 147, by contrast, the low-$\eta$ states are numerous (72--83) and spread
over $\theta$: the finite-size tail of the thermal distribution.

\subsection{$k=1$: exponential fragmentation; two pairs at the same $N$}
Rules 108, 156, 198 and 201 fire on a single pattern (108 on $11$, 201 on $00$, 156 on $10$, 198 on $01$) and have
exponentially many Krylov sectors (R2, R19). They were run to $N=20$ (201 to $N=19$; see below). The class is compared
at the largest $N$ that all four rules share, $N=19$ (Table~\ref{tab:summary}), and not each rule at its own maximum.
\begin{itemize}
\item \textbf{108 and 201} (fire on $11$ or $00$; mirror symmetric) have large blocks: the largest has 2583 and 5544
  states at $N=19$. Pooled over them, $\langle r\rangle=0.3922\pm0.0010$ (108) and $0.3891\pm0.0009$ (201), with
  $7\times10^4$ to $10^5$ ratios. This is Poisson ($0.3863$) up to a small, statistically resolved excess of
  0.003 to 0.006, which is stable from $N=17$ on.
\item \textbf{156 and 198} (fire on $10$ or $01$; a reflection pair, computed independently and identical) stay
  \emph{tiny}: the largest block has 100 states at $N=19$ (125 at $N=20$), and only $2139$ of the $2^{19}$ quasienergies are
  distinct. $\langle r\rangle=0.4209\pm0.0038$ at $N=19$ and $0.4172\pm0.0023$ at $N=20$, stable from $N=17$ on.
  This is clearly \emph{above} Poisson but far below GOE.
\end{itemize}
The block structure explains both values. The commutant study (R-T18/R-T27) shows that a block is a tensor product of
independent rooms, so its quasienergies are sums of room quasienergies. When many rooms of different sizes contribute
(108, 201), the sums are effectively independent and the statistics are Poisson. When a block is built from a few small,
repeated rooms (156, 198), the sums are \emph{structured}. Exact coincidences are removed, but (our reading, not tested) the surviving levels keep
an arithmetic regularity that suppresses small spacings. Neither pair shows level repulsion growing with $N$.

The gate singlets are exponentially many (at $N=19$: 52\,688 for 108, 37\,812 for 156 and 198, and 39\,872 for 201;
roughly doubling with each added site). These are the dark states of the rooms.

\subsection{$k=2$, XOR: arithmetic spectra}
150 and 105 have only a handful of Krylov sectors (the domain-wall shells, R8) and large blocks, but their spectra are
\emph{arithmetic}. The number of distinct quasienergies of 150 is exactly $3^m$ for even $N$ and $(3^m+(-1)^m)/2$ for odd
$N$ ($m=\lceil N/2\rceil$) at every $N=6\ldots14$, reproducing R8's law from an independent diagonalisation. 105 shares
the even-$N$ counts but has far fewer at odd $N$ (251 against 1093 at $N=13$). Up to 95\,\% of the eigenstates lie in
degenerate eigenspaces, and spacings repeat exactly. $\langle r\rangle$ (0.42 and 0.36) is therefore no measure of chaos
here. The right description is the free wall walk with sign defects of R8. Each Krylov sector contains one gate singlet at
$\theta=0$, an entangled dark state ($S$ between 0.13 and 1.14).

\subsection{$k=2$, one-sided: a boundary-driven tower, statistics undecided}
In 60/102 (gate fires iff the left or right neighbour is 1) and 153/195 (iff it is 0), the gate on a site is controlled by
one neighbour only. With open boundaries the gate algebra then produces a tower:
\begin{itemize}
\item 153/195: $N+1$ blocks of sizes $2^{N-1},2^{N-2},\ldots,2,1$, plus two singlets;
\item 60/102: $2N-4$ singlets and many small blocks.
\end{itemize}
At $N=5$ the dominant basis states of block $j$ (153) are those whose first site from the left in the $-1$ eigenstate of
$H$ is site $j$. The sites to its left are frozen in the dressed $+1$ state, and the right part is a smaller copy of the
chain. The entanglement accordingly falls into bands, one per block (Fig.~\ref{fig:rule153}). Within the blocks,
$\langle r\rangle$ is near Poisson but \emph{increases} steadily with $N$: for 60, $0.334, 0.351, 0.360, 0.369, 0.381,
0.395$ for $N=9\ldots14$; for 153, $0.37$--$0.414$ for $N=9\ldots13$. Integrability and a slow crossover towards GOE cannot
be told apart at these sizes \emph{on the open chain}. The tower itself is a boundary effect. On a ring it is absent, and
the same rules are chaotic with $\beta=1$ statistics ($\langle r\rangle=0.531$ for 153 at $N=14$; Sec.~\ref{sec:pbc}). So
on the open chain, the non-ergodicity of this class is produced by the ends.

\subsection{$k=0$ and $k=4$: trivial}
204 is the identity (one distinct quasienergy). In 51 every site gets $H$ in its layer, so $U=H^{\otimes N}$ exactly. Its
spectrum is $\{0,\pi\}$, every state is a gate singlet, and the eigenstates are products of single-site $H$ eigenstates,
with $S=0$.

\begin{figure}[p]
\centering
\ifgraphic{../../figures/r32/fig_r32_rule_H_W57.pdf}{\textwidth}
\caption{Rule 57 ($k=3$), largest $N$, obc0. (a)~Half-chain entanglement against quasienergy. Coloured: non-degenerate states,
by block. Grey: states in a degenerate eigenspace, whose entanglement depends on the basis. Diamonds: gate singlets (exact
eigenstates of every gate; product states at $\theta=0,\pi$). (b)~$\eta=S_2^\mathrm{Haar}/S_2^\mathrm{rand}(\mathrm{block})$,
basis independent. (c)~$P(r)$ of the fully resolved half spectrum against the Poisson, GOE and GUE surmises. (d)~$\langle
r\rangle$ against $N$ at three resolution levels.}
\label{fig:rule57}
\end{figure}

\begin{figure}[p]
\centering
\ifgraphic{../../figures/r32/fig_r32_rule_H_W147.pdf}{\textwidth}
\caption{Rule 147 ($k=3$), largest $N$; panels as in Fig.~\ref{fig:rule57}. In (d) the Krylov-only statistics (here the whole
spectrum) stay near Poisson, while resolving the spatial symmetry ($R$ for odd $N$, the half step $\Omega$ for even $N$)
brings $\langle r\rangle$ to the GOE value.}
\label{fig:rule147}
\end{figure}

\begin{figure}[p]
\centering
\ifgraphic{../../figures/r32/fig_r32_rule_H_W153.pdf}{\textwidth}
\caption{Rule 153 ($k=2$, one-sided), largest $N$, obc0; panels as in Fig.~\ref{fig:rule57}. The entanglement bands in (a) are the
blocks of the boundary-driven tower (sizes $2^{N-1},2^{N-2},\ldots$). $P(r)$ is close to Poisson, while $\langle r\rangle$
drifts upward with $N$.}
\label{fig:rule153}
\end{figure}
